% !TeX root = main.tex

\documentclass[a4paper,11pt,twoside]{report}

% Packages utiles
\usepackage[utf8]{inputenc}
\usepackage[Rejne]{fncychap} %Pour le style d'affichage des titres de chapitre
\usepackage[T1]{fontenc} % Encodage des fontes de caractères
\usepackage[french]{babel} % Utilisation des règles typographiques françaises
\usepackage[cyr]{aeguill} % Pour les guillemets et autres caractères
\usepackage[bitstream-charter]{mathdesign} % Police vectorielle charter
\usepackage{mathtools}
\usepackage{graphicx} % Insérer des images
\usepackage{parskip} % Sauter des lignes entre paragraphes
\usepackage{verbatim} % Intégration de code en type "Typewriter"
\usepackage{hyperref} % Liens hypertextes
\usepackage{emptypage} % Suppression du numéro de page, de l'entête et du pied de page pour les pages vides
\usepackage{fancyhdr} % Pour personnaliser l'entête et le pied de page
\usepackage{float}
\usepackage{url}
\usepackage{changepage}
\usepackage{lastpage}
\usepackage[justification=centering]{caption} % Pour centrer la légende des images
\usepackage{minitoc}

% Algo
\usepackage{algorithm}
\usepackage{algpseudocode}
\usepackage{import}
\usepackage{etoolbox}
\usepackage{multirow}
\usepackage{array}
\usepackage[caption=false]{subfig}

\graphicspath{{assets/}}

% Liste des algorithmes
\let\mylistof\listof
\renewcommand\listof[2]{\mylistof{algorithm}{Liste des algorithmes}}
% pour palier au problème de niveau des algos
\makeatletter
\providecommand*{\toclevel@algorithm}{0}
\makeatother

% Options de mise en page
\usepackage[tmargin=2.5cm,bmargin=2.5cm,lmargin=2.5cm,rmargin=2.5cm]{geometry} % Marges et format

% Métadonnées
\author{Travaillé Adrien}
\title{Projet de recherche technologique}
\date{\today}

% Définition de l'entête et du pied de page
\renewcommand{\headrulewidth}{1pt}
\renewcommand{\footrulewidth}{1pt}
\setlength{\headheight}{14pt}

% Nouvelle commande pour les tableaux
\newcolumntype{C}[1]{>{\centering\arraybackslash}m{#1}} %Utiliser C dans la définition des colonnes

% Largeur des traits de séparation dans les tableaux
\setlength{\arrayrulewidth}{0.5mm}

\begin{document}

\title{Projet de Fin d'Études \\[0.5cm] Développement d'un outil permettant le traitement et l'insertion des archives numériques sur Geofoncier au sein d'un cabinet de Géomètre-Expert}
\pagenumbering{gobble} %Ne pas numéroter la page de garde et de sommaire
\begin{titlepage}
  \centering
  {\Large INSA Strasbourg \par}
  \vspace{0.5cm}
  {\large Mémoire de Projet de Fin d'Études \par}

  \vspace{1.5cm}
  {\Huge \bfseries Projet de Fin d'Études \par}
  \vspace{0.8cm}
  {\Large \bfseries Développement d'un outil permettant le traitement et l'insertion des archives numériques sur Geofoncier au sein d'un cabinet de Géomètre-Expert\par}
  \vspace{0.2cm}
  {\Large \bfseries \par}

  \vspace{1.5cm}
  \begin{tabular}{ll}
    \textbf{Auteur} : & Travaillé Adrien \\
    \textbf{Encadrant} : & KOEHL Mathieu \\
    \textbf{Structure d'accueil} : & GEO-SIAPP \\
  \end{tabular}

  \vfill
  {\large Année universitaire : 2025-2026 \par}
\end{titlepage}



%------------------- Sommaire -------------------------------- %
\tableofcontents
\newpage

\chapter*{Remerciements}
\addcontentsline{toc}{chapter}{Remerciements}

% Remercier l’encadrant, l’équipe, l’établissement, les proches (1 page max).


\newpage

%%%%%%%%%%%%%%%%% Début de la numérotation des pages %%%%%%%%%%%%%%%%%%%%%%%%%%

\pagestyle{fancy}
\fancyhf{} % clear all fields
\fancyhead[R]{\rightmark}
\renewcommand{\sectionmark}[1]{%
\markright{\MakeUppercase{\thesection.\ #1}}}%
%\lhead{\leftmark}
%\lhead{}
\cfoot{\bfseries Page \thepage{} / \pageref{LastPage}}
\makeatletter
\let\ps@plain=\ps@fancy
\makeatother

%------------------- Inclusion des chapitres ----------------- %

\pagenumbering{arabic} % Commencer ici la numérotation
\setcounter{page}{1} % Forçage du compteur de page à 1
\chapter{Introduction}

\section{Contexte}
Ceci est un text en direct
\section{Problématique}
% Définir le problème et les contraintes.

\section{Objectifs}
\begin{itemize}
  \item Objectif 1.
  \item Objectif 2.
  \item Objectif 3.
\end{itemize}

\section{Organisation du document}
% 1–2 phrases par partie.



\chapter{Analyse du métier et des données}

\section{Le cadre de Geofoncier}
\subsection{Rôle et fonctionnement du portail}
% Décrire Geofoncier, à quoi ça sert, qui l’utilise, quels types de versements.

\subsection{Spécifications techniques des données attendues}
% Formats attendus, contraintes qualité, règles de nommage/tri, métadonnées, etc.

\section{Typologie des archives à traiter}
\subsection{Registres de communes et documents associés}
% Format, état, variations, structure des pages, présence de tableaux/colonnes.

\subsection{Plans cadastraux et pièces annexes (si pertinent)}
% À garder seulement si ton projet les traite réellement.

\subsection{Spécificités de l'écriture et difficultés}
% Manuscrit, bruit, encre, rotation, pliures, contraste, etc.


\chapter{État de l'art}

\section{Vision par ordinateur}
\subsection{Détection d'objets dans des documents}
% Pourquoi la détection/segmentation est utile pour isoler colonnes, cellules, zones texte.

\subsection{Pourquoi YOLO pour segmenter des documents}
% Arguments: robustesse, vitesse, écosystème, annotation, etc.

\section{Reconnaissance de l'écriture manuscrite (HTR)}
\subsection{RNN vs Transformers}
% Rappeler les grandes familles et leurs implications.

\subsection{Comparatif open-source (TrOCR, Kraken, EasyOCR)}
% Forces/faiblesses, contraintes d’entraînement, qualité sur manuscrit.

\section{Modèles Vision-Langage (VLM)}
\subsection{Apport d'un VLM pour lever les ambiguïtés}
% Exemple: arbitrage avec contexte visuel global.

\section{Synthèse}
% Ce que l’état de l’art ne résout pas entièrement et comment ton pipeline répond.


\chapter{Architecture du projet}

\section{Étape 1 : Segmentation et géométrie des pages}
\subsection{Détection des colonnes et cellules via YOLO}
% Données d'entraînement, classes, métriques, exemples.

\subsection{Reconstruction de la structure : interpolation \textit{Ghost Lines}}
% Décrire le besoin (traits manquants, défauts visuels) et la logique générale.

\section{Étape 2 : Moteur de transcription multi-modèles}
\subsection{Extraction des zones de texte}
% Comment les crops sont produits et normalisés.

\subsection{Stratégie multi-hypothèses}
% Générer plusieurs prédictions, score/confiance, sélection.

\section{Étape 3 : Arbitrage par VLM}
% Quand déclencher l’arbitrage et quel signal d’entrée/sortie.

\section{Protection des données sensibles}
% Mesures: anonymisation, stockage local, contrôle d’accès, journaux, etc.


\chapter{Fiabilisation et post-traitement des données}

\section{Décodage quand l'écriture est difficile}
% Contraintes de décodage, dictionnaires, caractères autorisés, etc.

\subsection{Contraintes au décodage (Logits Processors)}
% Expliquer l’idée: interdire hors-dictionnaire / hors commune, etc.

\section{Matching intelligent}
\subsection{Normalisation des toponymes}
% Accents, casse, abréviations, séparateurs, etc.

\subsection{Matching fuzzy et similarité visuelle}
% Levenshtein pondéré, confusions de lettres, Smith-Waterman, etc.

\section{Validation via des référentiels officiels}
% INSEE, listes officielles, contrôle de cohérence.


\chapter{Intégration sur Geofoncier et résultats}

\section{Structuration et export}
\subsection{Mise en forme des données extraites}
% CSV/Excel + sorties images de contrôle, traçabilité.

\subsection{Préparation des packages pour Geofoncier}
% Format final, organisation, règles de dépôt, vérifications.

\section{Évaluation des performances}
% Définir les métriques: taux de succès par commune, taux d’erreurs, etc.

\section{Analyse des limites}
% Faux positifs, besoin de contrôle humain, cas non couverts.



\chapter{Conclusion et perspectives}

\section{Synthèse des contributions}
% Contributions techniques: segmentation, HTR multi-modèles, arbitrage VLM, post-traitement.

\section{Apports personnels}
% Ce que tu as appris/produit: IA, développement, compréhension métier.

\section{Perspectives}
% Ouvertures: lecture d’autres archives, extension aux plans parcellaires, etc.



%------------------- Liste des figures et des tableaux ----------------------- %
\listoffigures
\thispagestyle{empty}
\listoftables
\thispagestyle{empty}
\listofalgorithms
\thispagestyle{empty}

%------------------- Bibliographie --------------------------- %
\newbool{firstbib}
\booltrue{firstbib}
\preto{\bibitem}{\ifbool{firstbib}{\thispagestyle{empty}\setbool{firstbib}{false}}{}}

% Configuration du style (BibTeX standard)
\bibliographystyle{apalike}

% Affiche toutes les références même celles non citées dans le texte
%\nocite{*}

% Affiche la bibliographie ici et indique le fichier .bib source
\bibliography{ref_prt}

\end{document}

